% figure-nested-intervals.tex — inline, no float wrapper.

\begin{center}
\begin{tikzpicture}[
  lbl/.style={font=\small},
  tick/.style={thin},
  arr/.style={->, >=stealth, thin}
]
\draw[thick] (0.5, 3.4) -- (9.5, 3.4);
\draw[tick, blue!70!black] (0.5, 3.2) -- (0.5, 3.6);
\draw[tick, blue!70!black] (9.5, 3.2) -- (9.5, 3.6);
\node[lbl, left]  at (0.4,  3.4) {$a_1$};
\node[lbl, right] at (9.6,  3.4) {$b_1$};
\node[lbl, left, blue!70!black] at (-0.5, 3.4) {$I_1$};
\draw[thick, blue!60!black] (1.8, 2.6) -- (7.8, 2.6);
\draw[tick, blue!60!black] (1.8, 2.4) -- (1.8, 2.8);
\draw[tick, blue!60!black] (7.8, 2.4) -- (7.8, 2.8);
\node[lbl, left]  at (1.7,  2.6) {$a_2$};
\node[lbl, right] at (7.9,  2.6) {$b_2$};
\node[lbl, left, blue!60!black] at (-0.5, 2.6) {$I_2$};
\draw[thick, blue!50!black] (3.0, 1.8) -- (6.2, 1.8);
\draw[tick, blue!50!black] (3.0, 1.6) -- (3.0, 2.0);
\draw[tick, blue!50!black] (6.2, 1.6) -- (6.2, 2.0);
\node[lbl, left]  at (2.9,  1.8) {$a_3$};
\node[lbl, right] at (6.3,  1.8) {$b_3$};
\node[lbl, left, blue!50!black] at (-0.5, 1.8) {$I_3$};
\draw[thick, blue!40!black] (3.8, 1.0) -- (5.2, 1.0);
\draw[tick, blue!40!black] (3.8, 0.8) -- (3.8, 1.2);
\draw[tick, blue!40!black] (5.2, 0.8) -- (5.2, 1.2);
\node[lbl, left]  at (3.7,  1.0) {$a_4$};
\node[lbl, right] at (5.3,  1.0) {$b_4$};
\node[lbl, left, blue!40!black] at (-0.5, 1.0) {$I_4$};
\draw[dashed, gray!60!black, thin] (4.5, 0.0) -- (4.5, 3.8);
\fill (4.5, -0.2) circle (3.5pt);
\node[lbl, below] at (4.5, -0.35) {$x = \sup\{a_n\}$};
\end{tikzpicture}
\end{center}
\captionof{figure}{%
  The Nested Interval Property (\cref{thm:nested-interval-property}).
  Left endpoints $a_n$ increase; right endpoints $b_n$ decrease; each
  $b_n$ is an upper bound for $\{a_n\}$. By the Axiom of Completeness,
  $x = \sup\{a_n\}$ exists and satisfies $a_n \leq x \leq b_n$ for all
  $n$, so $x \in \bigcap I_n$.%
}
\label{fig:nested-intervals}
\medskip
